% Section 4.2 (usage scenario: wildfire, temporal subsetting) -- the author's current draft,
% rendered as-is on 2026-11-07 (text unchanged; highlights, TODOs and comments kept).
% Every number is in ../CASE_STUDY_BRIEF.md.
%
% Preview with the real VGTC class (two columns, 9 pt):
%   TEXINPUTS=~/Desktop/unOrg/vis2026/vgtc-conference//: BSTINPUTS=~/Desktop/unOrg/vis2026/vgtc-conference//: \
%     latexmk -pdf -outdir=.texbuild section_4_2.tex
% To move the section into the paper, copy what lies between the "BEGIN SECTION" and
% "END SECTION" markers, and add refs_cs2.bib's entries to the paper's .bib.
\documentclass{vgtc}
\usepackage{times}
\renewcommand*\ttdefault{txtt}
\usepackage{mathptmx}
\usepackage{graphicx}
\usepackage{booktabs}
\usepackage{listings}
\usepackage{xcolor}
\usepackage{soul}
\vgtcinsertpkg
\graphicspath{{figures/}}

% ---- Preview only: the paper's macros ----
\newcommand{\sname}{Sieve}
\newcommand{\Sname}{Sieve}
\newcommand{\code}[1]{\texttt{#1}}
\newcommand{\todo}[1]{\textcolor{red}{\textbf{[TODO: #1]}}}
\lstdefinelanguage{Sieve}[]{Python}{morekeywords={source,fields,subsample,timesteps,region,save,render}}
\lstset{basicstyle=\ttfamily\scriptsize, breaklines=true, columns=fullflexible,
        frame=tb, captionpos=b, showstringspaces=false, aboveskip=4pt, belowskip=2pt}
\renewcommand{\lstlistingname}{Listing}
\definecolor{hlcolor}{RGB}{255,236,140}
\sethlcolor{hlcolor}
\soulregister\cite7 \soulregister\ref7 \soulregister\sname7 \soulregister\Sname7
\soulregister\textbf7 \soulregister\textit7 \soulregister\code1

\title{Section 4.2 preview}
\onlineid{0}
\vgtccategory{System}

\begin{document}

% ---- Preview only: the teaser (placeholder caption; the paper's teaser sits on page 1) ----
\begin{figure}[t]
  \centering
  \includegraphics[width=\columnwidth]{cs2_teaser}
  \caption{\todo{teaser caption} Teaser image. Inside the white box: the full-resolution
  pass; outside it: the coarse pass.}
  \label{fig:teaser}
\end{figure}

\setcounter{section}{3}
\section{Usage Scenarios}
\setcounter{subsection}{1}

% ===================== BEGIN SECTION (copy into the paper) =====================
\subsection{Tracking Wildfire Evolution with Time Windowing}
\label{sec:cs2}

\subsubsection{Data and Analysis Requirements}

The FIRETEC data \cite{linn2002firetec} simulates a fire spreading over a ridge. It is stored on a remote server as 88 timesteps of VTK files, each 1.1 GB (96.6 GB in total) holding nine variables on a $600\times500\times61$ terrain-following grid. Once the fire crosses the crest and moves down the far slope (the lee side), its spread slows down.

\paragraph{Scientific Task} In this scenario, the scientist asks whether at the point where the fire spread slows down significantly, the wind above the ground ahead of the fire reverses and blows against the spread. \todo{CITE where task is from} Answering this question requires two kinds of data. First, the scientist must locate \textit{when} the fire crosses the ridge and slows down. This involves sampling the series densely enough to identify the crossing and see where the speed drops; sampling density is not known in advance. Second, measuring the near-ground wind (a few meters above) requires the grid's full resolution, but only around the slowdown and near the fire.

% Without \sname{}, the scientist has three routes: (1) copy the whole series and meet both requirements, (2) a two-pass download by sampling the files followed by another copying of the chosen window, or (3) a script that reduces the data on the server. With \sname{}, both requirements get their own request, both issued from Paraview [cite] using the \sname{}'s plugin.

\hl{\sname{} allows the scientist to determine the relevant time window through exploration and then refine the data request around it. A spatially subsampled overview preserves all available timesteps to locate the slowdown; a subsequent request retrieves the selected window at full resolution to examine near-ground wind. Both requests are composed through \sname{}'s ParaView plugin [cite], combining \textbf{temporal-spatial-field} selections, and spatial subsampling as needed. This workflow avoids transferring the complete series or writing a custom server-side reduction script.}

\subsubsection{Coarse Pass: Locating the Slowdown}

The first request, a \emph{coarse pass}, reads every timestep, samples the grid by eight along each axis, and keeps only the five fields needed to follow the fire: the wind components (\code{u}, \code{v}, \code{w}), temperature (\code{theta}), and fuel density (\code{rhof\_1}).


\begin{lstlisting}[language=Sieve,numbers=none,basicstyle=\ttfamily\footnotesize, frame=tb, caption={\todo{Missing caption}}]
src = source("ssh://host//path/")
coarse = fields(src, ["rhof_1","theta","u","v","w"])
coarse = subsample(coarse, 8)
window = timesteps(src, [4000, 9000 - 45000])
window = region(fields(window, ["theta","u","v","w"]),
                x=(-100,200), y=(-120,120), z=(0,350),
                coords="world")
\end{lstlisting}

\begin{table*}[!h]
\centering\small
\begin{tabular*}{\textwidth}{@{\extracolsep{\fill}}lrrrrrr@{}}
\toprule
 & Timesteps & \multicolumn{3}{c}{GB moved} & Bytes vs & Script \\
\cmidrule(lr){3-5}
Route & surveyed & coarse pass & pass 2 & total & \Sname{} & lines \\
\midrule
\Sname{} & 88 & 0.05 & 0.36 & 0.41 & 1$\times$ & 6 \\
Two-pass download & 18 & 19.76 & 34.04 & 53.80 & 131$\times$ & 89 \\
Download everything & 88 & 96.62 & 0 & 96.62 & 234$\times$ & 73 \\
Scripts on the server & 88 & 0 & 0.61 & 0.61 & 1.5$\times$ & 72 \\
\bottomrule
\end{tabular*}
\caption{Bytes moved by each route to the same result. The first pass
surveys the series; the second fetches the window at full resolution.
\Sname{}'s requests were composed in ParaView, with no script.}
\label{tab:cs2}
\end{table*}

% The coarse pass fetched all 88 timesteps in 47\,MB. A two-pass download spends 19.8\,GB to sample every fifth file. Those 18 timesteps also show the front slowing after the crest, but they place the drop only within a 5,000-step interval (Figure~\ref{fig:cs2-speed}); the coarse pass resolves it to 1,000 steps for about 1/400 of the bytes. \hl{By reducing spatial detail while retaining every timestep, \sname{} lets the scientist identify the time window of interest \textit{without} committing in advance to a sampling interval that could obscure the onset of the slowdown.}

The coarse pass fetched all 88 timesteps in 47\,MB. Figure~\ref{fig:cs2-speed} shows front speed falling from about 20 to a few meters per 1,000 simulation steps as the fire crosses the crest and moves down the lee side. A two-pass download spends 19.8\,GB to sample every fifth file. \hl{The orange rings, calculated from the coarse-pass data at that spacing, show the broad slowdown but give speeds over 5,000-step intervals; the blue curve follows the transition at 1,000-step intervals. By subsampling space while retaining every timestep, \sname{} gives the scientist this closer view of the slowdown for about 1/400 of the transfer. The shaded band marks the main interval then selected for full-resolution inspection of near-ground wind.} By reducing spatial detail while retaining every timestep, \sname{} lets the scientist identify the time window of interest \textit{without} committing in advance to a sampling interval that could obscure the onset of the slowdown.


% the wind is close to zero at the lowest, so even small fluctuations in the air, flips the direction of wind. And in the data, it happens to turn slightly negative just as the fire crosses the crest, which makes it look like a wind reversal that would explain the slowdown.

% I need to rewrite it. "misleads" is not the correct term to use. Any scientist would look at the magnitude of the wind as well, so its not really misleading the scientist.

% [Huyen]: I rewrote that below

\subsubsection{Full-Resolution Pass: Examining Near-Ground Wind in the Selected Window}

The coarse pass shows a slightly negative wind component just as the fire crosses the crest. However, this occurs in the lowest grid layer, where wind speeds remain near zero and even small fluctuations can change the wind's direction. A sampling stride of eight retains the lowest grid layer and the next sampled layer at approximately \(12\,\mathrm{m}\) above ground, omitting the intermediate layers.

Using the time window identified in the coarse overview, \sname{} retrieved all 38 selected timesteps at full resolution, restricting the data to a box spanning the crest and lee slope and retaining only the wind components and temperature. A few meters above ground, the wind continues to blow in the direction of spread throughout the slowdown and reverses only about 16,000 simulation steps later. This later reversal therefore cannot explain the earlier slowdown. The analysis establishes \textit{when} the wind reverses, but not \textit{what} causes the reversal.


\begin{figure}[t]
\centering
\includegraphics[width=\columnwidth]{Case Study Diagrams/Wildfire_VTK/fig_speed_coverage.png}
\caption{Front speed over the run. Blue: \Sname{}'s coarse pass, every timestep
(0.05\,GB). Rings: every 5th file (19.8\,GB); each ring is the speed between two
consecutive files, plotted at their midpoint, from the coarse-pass data. Shaded:
the full-resolution window.}
\label{fig:cs2-speed}
\end{figure}
% ====================== END SECTION (copy into the paper) ======================

\bibliographystyle{abbrv-doi}
\bibliography{refs_cs2}

\end{document}
